\section{関連研究}
\label{sec:related}

\subsection{条件付き動作生成}

BCILに条件変数を導入する試みとして，Bi-LAT \cite{Kobayashi2025biLAT} がある．
Bi-LATは自然言語を条件として力プロファイルを変調した動作を生成でき，``softly'' や ``strongly'' のような指示で把持力を調整できる点で大きな進展である．
しかし言語によるラベルは定性的なカテゴリに留まり，``$3.5\,\text{Nm}$ で締め付ける'' のような定量的な数値指定ができない．また離散カテゴリ間の補間は困難で，新しい力レベルには再学習が必要になる．

ゴール条件付き模倣学習 \cite{Ding2019} は，ゴール画像やゴール位置ベクトルといった目標状態を条件として与え，そこへ到達する軌道を生成するアプローチである．
ゴールを数値で指定できる点は本研究と方向性が近いが，力情報を含む条件指定は対象外であり，``現在位置から $5\,\text{mm}$ 右へ'' のような操作量ベースの指定とも設計思想が異なる．

\subsection{残差ベースのロボット学習}

残差強化学習 \cite{Johannink2019} は，ベースポリシーに対する残差補正を強化学習で学習する手法であり，少数の補正を効率的に学習できる利点がある．
また，Davchev らの Residual Learning from Demonstration \cite{Davchev2022} は，動的動作プリミティブベースのポリシーに対して残差をデモンストレーションから学習することで，接触を含む操作タスクに適応させる手法を提案している．
さらに TER-DAgger \cite{Huang2026} は DAgger フレームワークに残差学習を組み込み，力認識フィードバックによって接触タスクの成功率を改善している．

これらの手法は残差という概念において本研究と方向性が共通しているが，
残差強化学習はオンラインの強化学習探索を，TER-DAgger はロールアウト中の人間介入を必要とするため，安全性の確保や事前に定めた数値目標への直接収束という点で課題がある．
本研究はオフラインの収録データのみを用い，かつ定量的なラベルと出力の間の単調増加性を活用して ILC による数値収束を目指す点で，既存の残差学習系手法と異なる．
